\documentclass[12pt,a4paper]{article}
\usepackage[utf8]{inputenc}
\usepackage[T1]{fontenc}
\usepackage[brazil]{babel}
\usepackage{amsmath,amssymb}
\usepackage{enumitem}
\usepackage{geometry}
\usepackage{fancyhdr}
\usepackage{listings}
\usepackage{xcolor}

\geometry{margin=2.5cm}
\pagestyle{fancy}
\fancyhf{}
\fancyhead[L]{GBC063 -- Inteligência Artificial -- Semana 2}
\fancyhead[R]{Prof. Marcelo Keese Albertini · FACOM · UFU}
\fancyfoot[C]{\thepage}
\renewcommand{\headrulewidth}{0.4pt}

\definecolor{codebg}{RGB}{32,38,46}
\definecolor{codefg}{RGB}{234,239,245}
\lstset{
  language=Python,
  basicstyle=\ttfamily\small,
  backgroundcolor=\color{codebg},
  keywordstyle=\color{blue},
  commentstyle=\color{gray},
  stringstyle=\color{green},
  frame=single,
  rulecolor=\color{gray},
  framerule=0.5pt,
  aboveskip=8pt, belowskip=8pt,
  showstringspaces=false,
  columns=flexible,
  keepspaces=true
}

\begin{document}

\section*{Lista de Exercícios -- Agentes, ambientes e NumPy}

\subsection*{Parte 1 -- PEAS e Propriedades do Ambiente}

\begin{enumerate}[label=\textbf{Q\arabic*.}, leftmargin=*]

\item Para cada um dos agentes abaixo, preencha a estrutura PEAS (Performance, Environment, Actuators, Sensors):

\begin{enumerate}
  \item Um robô aspirador de pó (2 cômodos)
  \item Um programa que joga xadrez
  \item Um carro autônomo
  \item Um sistema de recomendação de filmes (Netflix)
\end{enumerate}

\item Classifique cada um dos ambientes acima segundo as 6 propriedades (observável/parcialmente observável, determinístico/estocástico, episódico/sequencial, estático/dinâmico, discreto/contínuo, agente único/multiagente). Justifique cada classificação em uma frase.

\item O GridWorld 4×4 implementado em aula é um ambiente com as seguintes propriedades: totalmente observável, determinístico, sequencial, estático, discreto, agente único.

\begin{enumerate}
  \item O que mudaria na implementação de \texttt{step()} se o ambiente fosse \textbf{estocástico} (ex.: 20\% de chance de o movimento falhar e o agente não sair do lugar)?
  \item O que mudaria se fosse \textbf{parcialmente observável} (ex.: o agente só enxerga as células adjacentes)?
  \item O que mudaria se fosse \textbf{multiagente} (ex.: dois agentes competindo para alcançar o objetivo primeiro)?
\end{enumerate}

\subsection*{Parte 2 -- Tipos de Agente}

\item Descreva a diferença entre um agente \textbf{reativo simples} e um agente \textbf{com estado interno}. Dê um exemplo de uma situação em que o agente reativo falha e o agente com estado interno resolve.

\item Explique por que um agente \textbf{baseado em objetivo} precisa de \textbf{busca} para decidir. Qual é a relação entre o espaço de estados e a função \texttt{decidir()}?

\item O diagrama do \textbf{agente que aprende} (agente + crítico + elemento de aprendizado) é apresentado em aula como ``o diagrama do reinforcement learning''. Explique o papel de cada componente e como eles interagem.

\subsection*{Parte 3 -- NumPy e GridWorld}

\item Considere o seguinte código:

\begin{lstlisting}
import numpy as np

a = np.array([10, 20, 30, 40])
b = np.array([1, 2, 3, 4])
\end{lstlisting}

Qual é o resultado de cada operação abaixo? (Responda sem executar -- depois confira no Python.)

\begin{enumerate}
  \item \texttt{a + b}
  \item \texttt{a * b}
  \item \texttt{a > 25}
  \item \texttt{a[a > 25]}
  \item \texttt{np.argmax(a)}
  \item \texttt{np.sum(a * b)}
\end{enumerate}

\item O GridWorld usa \texttt{np.array} para representar a posição do agente. Por que usar um array NumPy em vez de uma tupla ou lista Python? Dê duas razões.

\item \textbf{Desafio de vetorização:} O código abaixo calcula a distância euclidiana entre um ponto e todos os pontos de uma lista usando laço \texttt{for}:

\begin{lstlisting}
import math

def distancias_loop(p, pontos):
    return [math.sqrt((p[0]-q[0])**2 + (p[1]-q[1])**2)
            for q in pontos]
\end{lstlisting}

Reescreva esta função usando NumPy vetorizado. Dica: use \texttt{np.linalg.norm} com o parâmetro \texttt{axis}.

\item \textbf{Implementação:} Modifique o GridWorld da aula para incluir:

\begin{enumerate}
  \item Um parâmetro \texttt{prob\_falha} no \texttt{\_\_init\_\_} que, se $> 0$, faz com que a ação escolhida tenha probabilidade \texttt{prob\_falha} de resultar em um movimento aleatório (ambiente estocástico).
  \item Um método \texttt{render()} que imprime a grade no terminal com caracteres: \texttt{.} para célula vazia, \texttt{\#} para obstáculo, \texttt{G} para objetivo, \texttt{A} para agente.
  \item Teste seu GridWorld estocástico com \texttt{prob\_falha=0.3} e execute o agente reativo 10 vezes. A taxa de sucesso (chegar ao objetivo) muda?
\end{enumerate}

\subsection*{Parte 4 -- Conexões}

\item O GridWorld que você implementou esta semana será reutilizado nas Semanas 10--14 para reinforcement learning. Explique: o que precisará mudar no código do GridWorld para que ele suporte Q-learning? O que \textbf{não} precisará mudar?

\item \textbf{Ponte para o Deep Learning:} O artigo ``Human-level control through deep reinforcement learning'' (Mnih et al., 2015, Nature) -- que introduziu o DQN -- usa a interface \texttt{reset()/step()} da biblioteca Arcade Learning Environment (ALE). Compare esta interface com a do GridWorld. O que há de fundamentalmente igual? O que há de diferente na escala?

\end{enumerate}

\end{document}
